\sect{श्रीपाद्मपुराणे माघमासमाहात्म्ये शिवरात्रिव्रतकथनम्}

\ganapatyadiDhyanam

\padmaPuranam

\dnsub{कथा}

\dnsub{अथ चत्वारिंशदधिकद्विशततमोऽध्यायः॥२४०}

\uvacha{दिलीप उवाच}

\threelineshloka
{केन व्याजेन वै व्याधो निराहारोऽभवन्मुने}
{कैलासं स कथं प्राप शिवरात्र्याश्च वैभवम्}
{श्रोतुमिच्छाम्यशेषेण वद विस्तरतो मुने} % 1

\uvacha{वसिष्ठ उवाच}

\twolineshloka
{शृणु राजन् प्रवक्ष्यामि शिवरात्रिव्रतं तव}
{व्रतानामुत्तमं साक्षाच्छिवलोकैकसाधनम्} % 2

\twolineshloka
{माघफाल्गुनयोर्मध्ये कृष्णपक्षे चतुर्दशी}
{शिवरात्रीति विज्ञेया सर्वपापापहारिणी} % 3

\twolineshloka
{कृतोपवासा ये तस्यां शिवमर्चन्ति जाग्रतः}
{बिल्वपत्रैश्चतुर्यामं ते यान्ति शिवतुल्यताम्} % 4

\twolineshloka
{न तपोभिर्न दानैश्च न वा जप्यसमाधिभिः}
{प्राप्यते तत्फलं राजन्नोपवासमखादिभिः} % 5

\twolineshloka
{गुह्याद् गुह्यतरं लोके व्रतमेतच्छिवप्रियम्}
{त्वयाऽपि खलु सर्वत्र न प्रकाश्यमिदं व्रतम्} % 6

\twolineshloka
{भूधराणां यथा मेरुस्तेजसां भास्करो यथा}
{द्विपदां च यथा विप्रः कपिलेव चतुष्पदाम्} % 7

\twolineshloka
{जप्यानामिव गायत्री रसानाममृतं यथा}
{पुरुषाणां यथा विष्णुः स्त्रीणां यद्वदरुन्धती} % 8

\twolineshloka
{शिवरात्रिव्रतं राजन् व्रतानामुत्तमं तथा}
{शिवरात्रिर्महावह्निर्भवानीशसमन्विता} % 9

\threelineshloka
{दहत्यवारितो योगाच्छुष्कार्द्रं कल्मषेन्धनम्}
{एतत्ते कथितं राजञ्शिवरात्रिव्रतं महत्}
{एवमेत्र पुरा देव्यै महादेवेन भाषितम्} % 10

\uvacha{दिलीप उवाच}

\onelineshloka
{कदा देव्या महादेवः कथं पृष्टस्तु तद् वद} % 11

\uvacha{वसिष्ठ उवाच}

\twolineshloka
{कैलासशिखरासीनं प्रसन्नमुखपङ्कजम्}
{त्रिलोचनं चतुर्बाहुं सर्वाभरणभूषितम्} % 12

\twolineshloka
{उमाधिष्ठितवामाङ्गं नागयज्ञोपवीतिनम्}
{वरदाभयहस्तं च नमज्जनवरप्रदम्} % 13

\twolineshloka
{व्याघ्रचर्मपरीधानं चन्द्रार्धकृतशेखरम्}
{गङ्गाप्लुतजटाजूटं भस्मगौरं वराननम्} % 14

\twolineshloka
{धारयन्तं महामालां ज्योतिश्चन्द्रार्कनिर्मलाम्}
{जगदुत्पत्तिसंहारस्थित्यनुग्रहकारिणम्} % 15

\twolineshloka
{महेष्वासमुदाराङ्गं कम्बुग्रीवं सुलोचनम्}
{सर्वाभरणसंयुक्तं शुक्लयज्ञोपवीतिनम्} % 16

\twolineshloka
{दृष्ट्वा प्रणम्य देवेशी प्रहर्षमतुलं गता}
{त्वरमाणाऽथ सङ्गम्य देवेशं वरवर्णिनी} % 17

\twolineshloka
{शुभां शुभावती चैत्र वदान्यां च सुमालिनीम्}
{समाहूयाऽऽगता देवी भूतस्त्रीभिश्च संयुता} % 18

\threelineshloka
{भवपादाब्जयुगले भक्त्या परमया युता}
{विकीर्य पुष्पजालानि सुमालिन्याऽऽहृतानि च}
{कृताञ्जलिपुटा भूत्वा पप्रच्छ शुभलोचना} % 19

\uvacha{देव्युवाच}

\twolineshloka
{अहोऽमृतमयी चैत्र कथा पापप्रणाशिनी}
{तत्र क्रीडाभिसम्बन्धात् त्वद्वाक्यामृतसम्भवा} % 20

\twolineshloka
{सुखावहसुखोद्गीणी दुःखक्षयविधायिनी}
{नीलोत्पलदलानां च मालेवोत्तरगन्धिनी} % 21

\twolineshloka
{नाद्यापि तृप्तिर्देवेश शृण्वत्या मम शङ्कर}
{प्रार्थितार्थान्यनेकानि दानं धर्मस्तथा परे} % 22

\twolineshloka
{यज्ञाचाऽऽयासबहुलास्तपांसि नियमानि च}
{बहूनि तानि लोकेऽस्मिन् पूतानि वृषकेतन} % 23

\twolineshloka
{व्रतानामुत्तमं देव भुक्तिमुक्तिप्रदायकम्}
{यन्न कस्यचिदाख्यातं त्वया सर्वार्थसाधकम्} % 24

\twolineshloka
{शीघ्रं निलीयते सर्वं पापं यच्छ्रवणादहो}
{तदहं श्रोतुमिच्छामि कथयस्व ममाग्रतः} % 25

\uvacha{महेश्वर उवाच}

\twolineshloka
{शृणु देवि व्रतं गुह्यं व्रतानामुत्तमोत्तमम्}
{यन कस्यचिदाख्यातं रहस्यं मुक्तिदायकम्} % 26

\twolineshloka
{येन वै कथ्यमानेनाप्यधर्मो विलयं व्रजेत्}
{तदहं कीर्तयिष्यामि शृणुष्वैकाग्रमानसा} % 27

\twolineshloka
{माघफाल्गुनयोर्मध्ये कृष्णपक्षे चतुर्दशी}
{शिवरात्रिरिति ख्याता सर्वयज्ञोत्तमोत्तमा} % 28

\twolineshloka
{दानयज्ञतपस्तीर्थव्रतकर्माणि यानि च}
{शिवरात्रिव्रतस्यापि कोट्यंशेन समानि न} % 29

\threelineshloka
{यैरिवं कलिहन्त्री च कृतान्तपथनाशिनी}
{भुक्तिदा मुक्तिदा देवि दिवानिशमुपोषिता}
{न ते यमपुरं यान्ति सत्यं सत्यं वरानने} % 30

\uvacha{देव्युवाच}

\twolineshloka
{कथं यमपुरं वन्ध्यं कथं शिवपुरं व्रजेत्}
{एतदेव महाश्चर्यं प्रत्ययं कुरु मे प्रभो} % 31


\uvacha{श्रीमहेश्वर उवाच}

\twolineshloka
{शृणु देवि यथा वृत्तं कथां पौराणिकीं प्रिये}
{कश्चिदासीत् पुराकाले निषादश्चाऽऽमिषप्रियः} % 32

\twolineshloka
{पर्वताग्रनिवासी च भूधरासन्नचारिभिः}
{समुत्पन्नैर्मृगैर्जीवन्कुटुम्ब परिपालकः} % 33

\twolineshloka
{आपीनांसो धनुष्पाणिः श्यामाङ्गः कृष्णकञ्चकः}
{बद्धगोधाङ्गुलित्राणो वामवाहौ च वर्मधृत्} % 34

\twolineshloka
{धनुर्वामे गृहीत्वा च दक्षिणे शरमुत्तमम्}
{निर्गतः स वनोद्देशान्निपादो मांसजीवकः} % 35

\twolineshloka
{वनं गतो निरीक्षिप्यन्सोऽन्तर्दिशमितस्ततः}
{वनमार्गे समन्विच्छन्नाश्रमे वनसूकरान्} % 36

\twolineshloka
{निराशो लुब्धकोऽतिष्ठद्यावदस्तं गतो रविः}
{चिन्तयञ्जनमासन्नं गतेऽर्के जीवघातक} % 37

\twolineshloka
{करिष्ये जागरं रात्रौ निश्चिता मम जीविका}
{गतोऽसौ जलमासनं तत्तीरे जालिमध्यतः} % 38

\twolineshloka
{प्रच्छन्नं कर्तुमारब्धवाऽऽत्मनो गुप्तिकारणम्}
{तत्र सन्तिष्ठते लिङ्गं स्वयम्भूतं वरानने} % 39

\twolineshloka
{सञ्छन्नं बिल्वविटपैः सपत्रैर्जालिमध्यतः}
{तानि विल्वस्य पत्राणि गृहीत्वा मार्गशोधने} % 40

\twolineshloka
{नीतानि दक्षिणे भागे न्यपतँल्लिङ्गमूर्धनि}
{न दिवा भोजनं तस्य ह्यामिषालुब्धचेतसः} % 41

\twolineshloka
{निरीक्षतः पुनस्तस्य न निद्राऽप्युपपद्यते}
{तस्य गन्धं समासाद्य लुब्धकस्य वरानने} % 42

\twolineshloka
{न तिष्ठन्ति मृगाः सर्वे शरघातभयात् तदा}
{तेन सा शर्वरी नीता ह्युदिते सूर्यमण्डले} % 43

\twolineshloka
{गतोऽसौ गृहमार्गेण निराशो धृततोमरः}
{मांसशून्यकरं दृष्ट्वा पिता पुत्रमभाषत} % 44

\uvacha{पितोवाच}

\onelineshloka
{नाऽऽनीतमामिषं पुत्र कथं रात्रिमुपोषितः} % 45

\uvacha{महेश्वर उवाच}

\onelineshloka
{ततो निषादः पितरं पृच्छन्तं प्रत्यभाषत} % 46

\uvacha{निषाद उवाच}

\twolineshloka
{नाssनीतमामिषं तात निराशाः शिशवो गताः}
{बाधते क्षुदपर्यन्ता स्वमितो भोजनं कुरु} % 47


\uvacha{महेश्वर उवाच}

\twolineshloka
{भोजनं तु कृतं तेन वृद्धेन सह भार्यया}
{धर्महीनो निषादस्तु धर्मवर्ती बनागतः} % 48

\threelineshloka
{अकामाज्जागराद्रात्रौ शिवरात्र्यां वरानने}
{मृतोऽसौ कालपर्यन्ते गृहीतो यमकिङ्करैः}
{शिवेन प्रेषितास्तस्मै विमानगणकोटयः} % 49

\uvacha{शिव उवाच}

\twolineshloka
{शीघ्रमानयत गत्वा प्राप्तो हि यमकिङ्करैः}
{निर्दग्धं किल्विषं तस्य शिवरात्र्यामुपोषणात्} % 50

\uvacha{महेश्वर उवाच}

\twolineshloka
{इति श्रुत्वा वचो दिव्यं गणास्ते गन्तुमुद्यताः}
{स्तुवन्तः परमं देवं शिवं शान्तमनामयम्} % 51

\twolineshloka
{गच्छन्तं च गणेशास्तमपश्यलुब्धकं तथा}
{गृहीतं यमदूतैस्तु पाशमुद्गरधारिभिः} % 52

\twolineshloka
{मुञ्चतैनं महात्मानमित्यूचुस्तान्गणेश्वराः}
{ततोऽब्रुवन्यमभटाः कस्मादेष विमुच्यताम्} % 53

\twolineshloka
{क्रूरासु यातना स्वेप पापिष्ठो जीवघातकः}
{ततस्ते हन्तुमारब्धाः खङ्गमुद्गरपट्टिशैः} % 54

\twolineshloka
{गृहीत्वैनं महात्मानं किङ्कराः कालचोदिताः}
{प्रावर्तत महायुद्धमन्योन्यवधकाङ्क्षिणाम्} % 55

\twolineshloka
{भिन्नमस्तककायाश्च शूलमुद्गरसायकैः}
{जर्जरीकृतदेहाच क्रन्दमानाः सुदारुणम्} % 56

\twolineshloka
{त्राहि त्राहीति गर्जन्तो गतास्ते यममन्दिरम्}
{निषादोऽथ गणैर्नीतो यत्र देवो महेश्वरः} % 57

\twolineshloka
{दृष्टमात्रः शिवेनासौ निषादः सुखतां गतः}
{ततोऽसौ दिव्यदेहस्थः कुण्डलाभरणोज्ज्वलः} % 58

\twolineshloka
{तस्मै दत्तं महेशेन विमानं सार्वगामिकम्}
{अक्षमालासमाकीर्ण पुष्पमालामलम्बितम्} % 59

\twolineshloka
{नानारत्नसमायुक्तं नानावादित्रघोषितम्}
{भुङ्क्ष्व भोगान्मम पुरे यावदाभूतसम्प्लवम्} % 60

\twolineshloka
{एवं लब्धवरो भूत्वा ह्यतिष्ठच्छिवशासने}
{अथ दूता वदन्त्यग्रे धर्मराजपुरस्थिताः} % 61

\twolineshloka
{रुधिरेणारुणाङ्गास्ते जर्जरीकृतमस्तकाः}
{प्रोचुश्च धर्मराजानं कृताञ्जलिपरिग्रहाः} % 62

\uvacha{दूता ऊचुः}

\twolineshloka
{शृणु राजन् यथावृत्तं युद्धं शिवगणैः सह}
{नीतोऽसौ पापकर्मा तु निषादो जीवघातकः} % 63

\twolineshloka
{अस्माभिर्नीयते राजन्कालमाप्तस्त्वदाज्ञया}
{एतस्मिन्नेव काले तु परमेशगणेश्वराः} % 64

\twolineshloka
{कालाग्निरुद्रसङ्काशाः शूलटङ्कगदाधराः}
{सिद्धाः सहस्रहस्ताथ त्रिनेत्राव जटाधराः} % 65

\twolineshloka
{दृष्टिनाः सर्वतोभद्रा भस्मपाण्डुरविग्रहाः}
{भुजङ्गहारवलयाः शशाङ्ककृतशेखराः} % 66

\twolineshloka
{गम्भीरोदण्डसंरावा. ब्रुवन्तश्च मुहुर्मुहुः}
{तत्राऽऽगत्य त्वरायुक्ताः प्रोचुरस्मानिदं वचः} % 67

\twolineshloka
{मुञ्चतैनं महात्मानं तपसा दग्धकिल्विषम्}
{श्रुत्वा गणेश्वरवाक्यमुक्तमस्माभिरप्ययम्} % 68

\twolineshloka
{न मोक्तव्यो निपादो हि पापात्मा जीवघातकः}
{अनेन घातिता जीवा असङ्ख्याता गणेश्वराः} % 69

\twolineshloka
{चित्राभियतनाभिस्तु बध्योऽयं हि यमाज्ञया}
{ते त्वस्मद्वचनं श्रुत्वा गणेशास्त्वतिगर्विताः} % 70

\twolineshloka
{शूलटङ्कगदाभिश्च खड्गमुद्गरतोमरैः}
{भिन्दिपालकुठारैश्च वज्रमुष्ट्युपलैस्तथा} % 71

\threelineshloka
{वयं हता गणैस्तैस्तैर्महाबलपराक्रमैः}
{वद्धस्तु विविधैः पाशैर्गृहीतो जीवघातकः}
{बहुनाऽत्र किमुक्तेन पुनस्तैरेव रक्षितः} % 72

\uvacha{महेश्वर उवाच}

\onelineshloka
{तैरेवमुक्तः सङ्क्रुद्धो धर्मात्मा जीवितेश्वरः} % 73

\uvacha{यम उवाच}

\twolineshloka
{पापिष्ठो जीवघाती च निपादो निर्गुगस्त्वयम्}
{कथं शिवपुरं याति चित्रगुप्त विचारय} % 74

\uvacha{चित्रगुप्त उवाच}

\threelineshloka
{निरीक्ष्य पुस्तकं तेन न किञ्चित्सुकृतं कृतम्}
{धर्मबुद्धिर्न तस्यास्ति धर्माधर्मौ\footnotemark{} न विन्दति}
{एतस्मिन्नखिलं ज्ञातं सत्यं सत्यं वदाम्यहम्} % 75

\footnotetext{धर्माधर्मं}

\uvacha{यम उवाच}

\threelineshloka
{इति ज्ञात्वा निषादस्य चित्रगुप्त निवेदितम्}
{सोऽहं चिन्तां करोमीह जन्तूनां पापकर्मणाम्}
{गत्वा निवेदयिष्यामि ह्ययोग्यं तैस्तु यत्कृतम्} % 76

\uvacha{महेश्वर उवाच}

\twolineshloka
{एवमुक्त्वा गतः शीघ्रं यत्राऽऽस्ते शङ्करः स्वयम्}
{दृष्ट्वा तु देवदेवेशं शङ्करं स्तोतुमैरयत्} % 77

\uvacha{यम उवाच}

\twolineshloka
{नमस्त्रैलोक्यनाथाय महाबलपिनाकिने}
{साक्षात्कालविनाशाय कालनिर्दाहिने नमः} % 78

\twolineshloka
{शिवागमार्णवान्तस्थज्ञानरत्नप्रदायिने}
{हृदि स्थिताय सर्वेषां साक्षिणे जगतां विभो} % 79

\twolineshloka
{अज्ञानतिमिरान्धस्य तमसोऽतीतमूर्तये}
{अनाश्रिताय तुष्टाय कपालाय नमोऽस्तु ते} % 80

\twolineshloka
{अनादिमलभेत्रे च चिगुणोदयहेतवे}
{गुणप्रदाय गूढार्थद्योतकाय नमोऽस्तु ते} % 81

\uvacha{महेश्वर उवाच}

\onelineshloka
{एवं स्तुत्वा महादेवं प्रणामो दण्डवत्कृतः} % 82

\uvacha{यम उवाच}

\twolineshloka
{मेरुमन्दरतुल्यानि दुष्कृतानि बहून्यपि}
{नश्यन्ति तानि सर्वाणि तव पादाब्जचिन्तया} % 83

\uvacha{महेश्वर उवाच}

\onelineshloka
{इत्युक्त्वा दण्डमुद्रा तु पादाब्जे तु निवेदिता} % 84

\uvacha{शिव उवाच}

\twolineshloka
{किमर्थं दण्डमुद्राऽपि त्यक्तेयं यम सत्तम}
{केनापराधिना धर्मो धर्मराज उपेक्षितः} % 85

\uvacha{धर्मराज उवाच}

\twolineshloka
{त्वद्गणैर्देवदेवेश जगतां पालन प्रभो}
{मदीयाः किङ्करा देव घातिताः शक्तिमुद्गरैः} % 86

\twolineshloka
{निषादो जीवघाती च सर्वकर्मवहिष्कृतः}
{मांसलुब्धश्च देवेश विचचार महावने} % 87

\twolineshloka
{न लब्धं पिशितं तेन निपादेन धनुष्मता}
{जलासन्नगतो रात्र्यां गतेऽर्के जीवघातकः} % 88

\twolineshloka
{मृगा वा न भवन्त्यस्य न निद्रा च भवत्यहो}
{न कृतं भोजनं दिश्र्या शिवरात्र्यामुपोषितः} % 89

\twolineshloka
{अपर्यन्तां क्षुत्रां प्राप्य उदिते सूर्यमण्डले}
{गृहं गतोऽसौ देवेश त्वसम्प्राप्तमनोरथः} % 90

\twolineshloka
{पापमेवाकरोत्पापी निषादो मांसविक्रयी}
{अनेन सुकृतं देव न किञ्चिदुपपादितम्} % 91

\twolineshloka
{विचित्रयातनार्हस्तु पापिष्ठो जीवघातकः}
{सर्वाल्लोकान्विनिर्जित्य गणेश्वरमवाप्तवान्} % 92

\twolineshloka
{देवदेव महादेव भक्तानामार्तिनाशन}
{किं कृत्यमत्र देवेश त्वदाज्ञाकारिणा मया} % 93

\uvacha{महेश्वर उवाच}

\twolineshloka
{इत्युक्तो धर्मराजेन भगवान्भक्तवत्सलः}
{प्राह गम्भीरया वाचा शिवरात्रिमनुस्मरन्} % 94

\uvacha{शिव उवाच}

\twolineshloka
{अयं हि शुद्धः शबरः पुण्यात्मा धार्मिको महान्}
{तपस्वी मत्प्रियो नित्यां शिवरात्रिमुपोषितः} % 95

\twolineshloka
{शिवरात्रिरिति ख्याता माघकृष्णचतुर्दशी}
{भुक्तिमुक्तिप्रदा नित्यं सर्वपापप्रणाशिनी} % 96

\twolineshloka
{एवं मङ्गलदाऽभीष्टप्रदा पुण्यविवर्धिनी}
{यमशासनहत्री च श्रीपदायोगदायिनी} % 97

\twolineshloka
{सम्यक्सिद्धिकरी पूज्या सौभाग्यफलदायिनी}
{निर्मितं हि मया पूर्व सुजागरमनुत्तमम्} % 98

\threelineshloka
{व्रतं तस्यां तिथौ शैवं लोकानां हितकाम्यया}
{शिवरात्रिप्रभावेन कृतार्थः शवरः स्वयम्}
{जीवितेश वरं ब्रूहि वरदोऽहं तवेप्सितम्} % 99

\uvacha{यम उवाच}

\twolineshloka
{कृपालय महादेव भक्तानामभयप्रद}
{संसारसागरभ्रान्तिपरिविच्छेदिने नमः} % 100

\twolineshloka
{नमः पिनाकहस्ताय नमस्ते कृत्तिवाससे}
{तव पादाब्जयुगले भक्ति देहि महेश्वर} % 101

\uvacha{महेश्वर उवाच}

\onelineshloka
{भवत्वित्याह भगवान्गच्छ त्वं नगरं प्रति} % 102

\fourlineindentedshloka
{इत्येवमुक्तो वृषकेतनेन}
{प्रहृष्टरोमावृतसर्वगात्रः}
{व्रतप्रभावं शबरेण लब्धं}
{पश्यन् पदं स्वं भवनं जगाम} % 103


॥इति श्रीमहापुराणे पाद्म उत्तरखण्डे माघ माहात्म्ये वसिष्टदिलीपसंवादे शिवरात्रि प्रभावकथनं नाम चत्वारिंशदधिकद्विशततमोऽध्यायः॥२४०॥


\dnsub{अथैकचत्वारिंशदधिकद्विशततमोऽध्यायः॥२४१}
\resetShloka

\uvacha{देव्युवाच}

\twolineshloka
{श्रुतो व्रतानुभावश्च त्वन्मुखाम्भोरुहान्मया}
{चरितं शवरस्यापि जीवितेशगतिस्तथा} % 1

\twolineshloka
{नानार्थदं महादेव श्रुत्वा वाक्यामृतं च ते}
{प्रीतिरस्य स्वरूपं यत्पुनर्मे वक्तुमर्हसि} % 2

\uvacha{महेश्वर उवाच}

\twolineshloka
{जीवघाती च शबरः शिवरात्र्यामुपोषितः}
{अबुद्धिपूर्व देवेशि गाणपत्यमवाप्तवान्} % 3

\twolineshloka
{श्रद्धयाऽभीप्सया प्रीत्या भीत्या च हृदयेन वा}
{कृत्वा च जागरं रात्रौ मुच्यते सर्वकिल्बिषैः} % 4

\twolineshloka
{प्रभातसमये बुद्ध्वा गुरोर्गेहं समागतः}
{तस्याऽऽज्ञां प्रार्थयेत्पूर्वं व्रतानुचरणाय वै} % 5

\twolineshloka
{स्नात्वा शुक्लाम्बरधरः कृतमौनो जितेन्द्रियः}
{कृताहिकविधिस्तत्र गते चार्के समाहितः} % 6

\twolineshloka
{स्नात्वा वैनायकीं पूजां कल्पयेत्पुरतः शुभैः}
{यजनं प्रतियामं च पाद्याघैरागमोदितैः} % 7

\twolineshloka
{कृताभ्यङ्गाभिषेकं च प्रतियामं समाचरेत्}
{पञ्चगव्यादिभिश्चैव नालिकेरफलोदकैः} % 8

\twolineshloka
{अथो\footnotemark{}ऽन्यैरभिषेकाहैरोषधीविल्वपत्रकैः}
{स्नापयेच्च महादेवं सहस्राद्यैश्व शाम्भवैः} % 9

\twolineshloka
{पिष्टामलकहारिद्रचूर्णैरुद्वर्तयेत्सुधीः}
{अर्चयेद्विल्वपत्रैश्च गन्धतोयैश्व सेवयेत्} % 10

\twolineshloka
{स्वर्णोदकै रत्नतोयैश्चाभिषेकं समाचरेत्}
{तान्तवेनाथ निर्मृज्य नीराजनमथाऽऽचरेत्} % 11

\twolineshloka
{वस्त्रैर्नानाविधैश्चैव विशेषैर्धृपितैस्तथा}
{संवेष्टयेयथाशोभं सौवर्णैर्भूषणैरपि} % 12

\twolineshloka
{अलङ्कृत्य महादेवं पूजयेद्विल्वपत्रकैः}
{जातीचम्पकपुन्नागपद्मोत्पलकदम्बकैः} % 13

\twolineshloka
{कर्णिकारनवश्वेतमन्दारकुरवैस्तथा}
{मल्लिकाशोकधत्तूरशम्यर्कारग्वधैस्तथा} % 14

\twolineshloka
{करवीरयवाङ्कोलनन्द्यावर्तपलाशकैः}
{तुलसीनागकोरण्टकुसुमैश्च सुपूजयेत्} % 15

\twolineshloka
{नीलोत्पलैर्विशेषेण पूजयेल्लिङ्गमेश्वरम्}
{धूपदीम्पैश्च नैवेद्यैस्ताम्बूलघृतदीपकैः} % 16

\twolineshloka
{अशेषैर्वहुभक्षैव भोज्यैव विविधैरपि}
{जागरगीतनृत्याद्यैः प्रदीपाद्युपहारकैः} % 17

\twolineshloka
{सूर्यघाम्पैरनेकैव वीणावेणुरवैस्तथा}
{स्तोत्रमङ्गलवाद्यैश्च वेदघोषैरनेकशः} % 18

\twolineshloka
{शिवधर्मपुराणाद्यैर्मत्र महेश्वरोक्तकैः}
{प्रदक्षिणनमस्कारप्रणवैश्चुलुकोदकैः} % 19

\footnotetext{सन्धिरार्षः}

\twolineshloka
{एवं\footnotemark{} प्रजागरं कुर्यात्प्रतियामं विशेषतः}
{ध्यानं च श्रवणं नित्यं शिवधर्मागमेन तु} % 20

\footnotetext{रात्रौ}

\twolineshloka
{शिवमन्त्रजपं कृत्वा प्रभाते विमले पुनः }
{दानं भक्त्या\footnotemark{} च भक्तानां ब्राह्मणानां विशेषतः} % 21

\footnotetext{दत्त्वा}

\twolineshloka
{बालवृद्धातुराणां च शक्त्या च परितोषयेत्}
{योगिनामन्नपानाभ्यामक्षसूत्रकमण्डलू} % 22

\twolineshloka
{कौपीनाच्छादनं दण्डं भिक्षापात्रं च भस्म च}
{ददद्वित्तानुसारेण गुरुं सम्पूजयेत्ततः} % 23

\twolineshloka
{हेमाङ्गुलीयवस्त्राद्यैर्गन्धपूजा\footnotemark{}दिभिः शुभैः}
{सम्पूज्य प्रार्थयेत्पश्चात्कृताञ्जलिपरिग्रहः} % 24

\footnotetext{पुष्पा}

\twolineshloka
{त्वदाज्ञया कृतं सर्व शिवरात्रिमहाव्रतम्}
{अनुगृह्णीष्व मां नित्यमपराधं क्षमस्व मे} % 25

\twolineshloka
{व्रतं कर्तुमशक्तश्चेश्चलचित्तश्च मानवः}
{अथवाऽन्यप्रकारेण जागरं कारयेत्पुनः} % 26

\twolineshloka
{गेयनृत्योपहारैश्व स्तोत्रमङ्गलवादनैः}
{नानाश्चर्यप्रदानैर्वा नानाविधफलान्वितैः} % 27

\twolineshloka
{शिवरात्र्यां विशेषैर्वा शिवक्षेत्रे विशेषतः}
{येन केनाप्युपायेन शिवरात्र्यामुपोषितः} % 28

\twolineshloka
{जागरं कारयेद्धीमान्पातकैः स प्रमुच्यते}
{पुरुषो वाऽथ नारी वा कुमारो वाऽथ कन्यका} % 29

\twolineshloka
{किं न विन्देत देवेशि शिवरात्रिमहोत्सवे}
{भर्तृहीना च या नारी शिवराज्युपवासतः} % 30

\twolineshloka
{तस्मातिष्ठति मङ्गल्यं सुचिरं कालमक्षयम्}
{चान्द्रायणसहस्रैश्च प्राजापत्यशतैरपि} % 31

\twolineshloka
{मासोपवासैरन्यैश्च यत्फलं लभते च सः}
{ततः कोटिफलं\footnotemark{} लब्धं शिवरात्रिप्रजागरात्} % 32

\footnotetext{गुणं}

\twolineshloka
{सर्वयज्ञतपोदानतीर्थवेदेषु यत्फलम्}
{तत्सर्वं लभते देवि शिवरात्रिमहाव्रती} % 33

\twolineshloka
{संवत्सरं प्रतिदिनं तपसा यत्फलं भवेत्}
{तत्सर्वं त्रिगुणीकृत्य शिवरात्र्यामुपोषितः} % 34

\twolineshloka
{जन्मकोटिसहस्रैस्तु यत्फलं पूर्वसञ्चितम्}
{तत्फलं तस्य देवेशि शिवरात्रिप्रजागरात्} % 35

\twolineshloka
{ब्रह्महा गुरुघाती च वीरहा भ्रूणहा तथा}
{मयपश्च तथा गोनो मातृहा पितृहा तथा} % 36

\twolineshloka
{स्तेयी सुवर्णस्तेयी च गुरुतल्परतः सदा}
{मुच्यते नृपलीसक्तः शिवरात्रिप्रजागरात्} % 37

\twolineshloka
{परदारप्रधर्षी च देवब्रह्मस्वहा तथा}
{मुच्यते मित्रघाती च कृतन्नोऽपि वरानने} % 38

\twolineshloka
{विव\footnotemark{}प्रच्यावकश्चैव लिङ्गप्रध्वंसकस्तथा}
{मुच्यते नात्र सन्देहः शिवरात्र्यां शिवार्चनात्} % 39

\footnotetext{वीर्य}

\twolineshloka
{वाचिकानि विचित्राणि मानसानि महान्ति च}
{कायिकानि सहस्राणि तथा सांसर्गिकाणि च} % 40

\twolineshloka
{भित्त्वा विमुच्यते सर्वः शिवरात्रिप्रजागरात्}
{अस्थिमज्जागतं पापं सर्वजन्मान्तरैरपि} % 41

\twolineshloka
{बुद्ध्याऽबुद्ध्या च देवेशि यदि वा वारुणीं पिवेत्}
{मुच्यते नात्र सन्देहः शिवरात्रिप्रजागरात्} % 42

\twolineshloka
{अजपित्वा हुताशी च अदाता च विमुच्यते}
{यो लब्ध्वा देवि मानुष्यमल्पस्वपि च जन्मसु} % 43

\twolineshloka
{अर्चयेदैश्वरं लिङ्गं विघ्नेशं षण्मुखं तथा}
{अभीत्य शिवविद्यां च परेभ्यो न वदन्ति ये} % 44

\twolineshloka
{विवृण्वन्ति न शृण्वन्ति तमोपहतचेतसः}
{तेषां पापानि नश्यन्ति शिवरात्रिप्रजागरात्} % 45

\twolineshloka
{ये निन्दन्त्यैश्वरं मार्गमाश्वर्यं धर्मदर्शनम्}
{वेदांश्च शिवभक्तांश्च वैदिकाचारमेव वा} % 46

\twolineshloka
{नश्यन्ति तानि\footnotemark{} पापानि शिवरात्रिप्रजागरात्}
{अर्चितं शङ्करं दृष्ट्वा न नमन्त्यल्पबुद्धयः} % 47

\footnotetext{तेषां}

\twolineshloka
{येषां न राजते देवि ललाटं भस्मकैः शुभैः}
{तेषां पापानि नश्यन्ति शिवरात्रिप्रजागरात्} % 48

\twolineshloka
{उत्तमाङ्गे जटा येषां संसारभयनाशिनी}
{प्राञ्जलि वा सदा मह्यं न नमन्ति च शोभने} % 49

\twolineshloka
{तेषां पापानि नश्यन्ति शिवरात्रिप्रजागरात्}
{न पश्यन्त्यैश्वरं लिङ्गं दिनं प्रत्यमरेश्वरि} % 50

\twolineshloka
{ये तु वा नापि गच्छन्ति शिवक्षेत्रेषु मानवाः}
{तेषां पापानि नश्यन्ति शिवरात्रिप्रजागरात्} % 51

\twolineshloka
{ये च ब्रह्मादिभिस्तुल्यं त्वां मां लक्ष्म्या च शक्तिभिः}
{गुरुं ये प्राकृतैः सार्धं संस्मरन्ति वदन्ति ये} % 52

\twolineshloka
{तेषां पापानि नश्यन्ति शिवरात्रिप्रजागरात्}
{पर्वमैथुन कर्तारः परदाराभिगामिनः} % 53

\twolineshloka
{ये परित्यागसंयुक्ताः पुनः सङ्केन वाचिताः}
{शिवलिङ्गं महापुण्यं ये स्पृशन्ति न ते कचित्} % 54

\twolineshloka
{जन्मस्वनेकेषु देवि सादरं धियते न वा}
{तेषां पापानि नश्यन्ति शिवरात्रिप्रजागरात्} % 55

\uvacha{वसिष्ठ उवाच}

\twolineshloka
{इत्येवं कथितं विप्रा\footnotemark{} देव्यै देवेन भाषितम्}
{शिवरात्र्याश्र माहात्म्यं देवानां सन्निधौ पुरा} % 56

\footnotetext{राजन्}

\twolineshloka
{मत्पूर्वगणपाः सर्वे देवा ब्रह्मपुरःसराः}
{मुनयश्च महात्मानः सनातनपुरोगमाः} % 57

\twolineshloka
{कैलासवासिनः सर्वे मेरौ देवनिकेतने}
{प्रीता बभूवुर्विप्रेन्द्राः श्रद्धां कृत्वा कुतूहलात्} % 58

\twolineshloka
{तदाप्रभृति ब्रह्माद्याः शिवरात्रिमहाव्रतम्}
{कुर्वन्ति गौरवात् सर्वे शिवस्य परमात्मनः} % 59

\twolineshloka
{तस्मात् सर्वप्रयत्नेन कर्तव्यं तच्छिवाज्ञया}
{शिवरात्रिव्रतं विप्राः\footnotemark{}  कीर्तयेद्यः शृणोति वा} % 60

\footnotetext{राजन्}

\twolineshloka
{सर्वपापविनिर्मुक्तः शिवलोके महीयते}
{यत्रेदं कीर्त्यते विप्र\footnotemark{}  देव्या सन्निहितः शिवः} % 61

\footnotetext{भूप}

\twolineshloka
{तत्र देवाः सगन्धर्वाः सर्किनरमहोरगाः}
{तिष्ठन्त्यनुग्रहं कर्तुं शिवरात्र्यां मजागरात्} % 62

\twolineshloka
{अस्याध्यायस्य योऽर्थज्ञः स विज्ञेयः सदाशिवः}
{तं पूजयेन्महात्मानं भुक्तिमुक्तिप्रवर्तकम्} % 63

\twolineshloka
{महापातकयुक्तो वा युक्तो वा सर्वपातकैः}
{दोषैः कृतैर्न लिप्येत व्रतानुश्रवणादहो} % 64

\twolineshloka
{कीर्तनीयमिदं सद्भिः शिवरात्र्यां विशेषतः}
{सर्वतीर्थेषु यत्पुण्यं सर्वयज्ञेषु यत्फलम्} % 65

\twolineshloka
{तत्सर्व कोटिगुणितं प्राप्नोति श्रवणादपि}
{यत्पुण्यं शिवरात्र्यां वै व्याख्यानेनैव तद्भवेत्} % 66

\twolineshloka
{यान्यान्प्रार्थयते कामांस्तांस्तान्प्राप्नोति मानवः}
{व्रतानुश्रवणादेव तत्पुण्यं लभते नरः} % 67

\uvacha{दिलीप उवाच}

\twolineshloka
{ऐहिकामुष्मिकाभीष्टप्रदं परमदुर्लभम्}
{चतुर्थनिर्वृतेर्मार्गं चतुर्वर्गफलप्रदम्} % 68

\twolineshloka
{शैवं व्रतं सदा शुद्धं विशेषात्पुत्रदायकम्}
{तपश्चर्यापुरस्कारम् अस्माकं वद सूतज\footnotemark{}} % 69

\footnotetext{तन्मे वद महामुने}

\uvacha{सूत (वसिष्ठ) उवाच}

\twolineshloka
{व्रतानामुत्तमं विप्रा\footnotemark{} विशेषात्पुत्रदायकम्}
{ऐहिकामुष्मिकाभीष्टमदं परमदुर्लभम्} % 70

\footnotetext{भूप}

\twolineshloka
{शृणुध्वं\footnotemark{} सोमनाथस्य सौराष्ट्रपुरवासिनः}
{शम्भोरुमासहायस्य सर्वमङ्गलदायिनः} % 71

\footnotetext{शृणु त्वं}

\twolineshloka
{सौराष्ट्रमिति विख्यातं नाम्ना त्रैलोक्यपूजितम्}
{पुत्रकामो बहून्पुत्रान्धनमिच्छुर्धनं लभेत्} % 72

\twolineshloka
{वन्ध्याऽपि लभते पुत्रान्कन्या सत्पतिमाप्नुयात्}
{कन्यार्थी लभते कन्यां जयकामो जयं लभेत्} % 73

\twolineshloka
{शत्रुभिर्नाभिभूयेत पुत्रपौत्रैश्च संयुतः}
{पठमानो लभेत्सर्व स्वयमुद्दिश्य पठ्यते} % 74

\twolineshloka
{नाकाले मरणं तस्य न सर्वैर्दृश्यते ततः}
{न विषं क्रमते देहे जडान्धत्वं न मूकता} % 75

\twolineshloka
{न चोत्पातभयं तस्य वृश्चिकस्य भयं तथा}
{नाभिचारकृतैर्दोषैर्न लिप्येत कदाचन} % 76

\twolineshloka
{यत्रेदं लिखितं विप्राः\footnotemark{} स्थापितं पूजितं यदि}
{न तत्र चैकं दौर्गत्यं कदाचिदपि जायते} % 77

\footnotetext{भूप}

इति श्रीमहापुराणे पाद्म उत्तरखण्डे माघमाहात्म्ये वसिष्ठदिलीपसंवादे शिवरात्रिव्रताख्यानं नामैकचत्वारिंशदधिकद्विशततमोऽध्यायः॥२४१॥


\genMangalaShloka